% ===== §13 The cultivation of the field and the co-creation of value =====
\section{The Cultivation of the Field and the Co-Creation of Value}
\label{p09:sec:cocreation}

\noindent
Perception clears the ground; it does not yet plant. Having learned to perceive the quiet third value against the grain of the loud two (\xref{p09:sec:perception}), the practical sequence can take its next step: from seeing what is generated to \emph{fostering its generation}. This is two questions that turn out to be one---how the field is cultivated (the lower half of the question \xref{p09:sec:eudaimonia} left open) and how value is co-created---and the burden of this section is to show why they are one, and what, concretely, the cultivating and the co-creating consist in.

They are one question because, on the account of \xref{p09:sec:eudaimonia}, the field just is the condition under which each participant's own dynamics unfold and generate value, and co-creation just is that unfolding. To cultivate the field is to make co-creation possible; co-creation is what a cultivated field is for. The seeming circularity is not a defect but the shape of the matter: \emb{one cultivates a field by co-creating in it, and one co-creates by cultivating the field for the others.} The two are the same activity described from the side of the condition and from the side of the act.

\subsection{What cultivation is not, and therefore is}
\label{p09:sec:cultivation-negative}

The surest approach to what cultivation \emph{is} runs through what it is not, for the whole of the paper's negative wisdom converges here, and cultivation is most easily understood as the positive name for a set of refusals already established.

To cultivate a field is \emph{not} to present a structure (\xref{p09:sec:eudaimonia}): not to lay down roles, scripts, an arrangement of who does what. It is \emph{not} to prescribe the path the flow must take, for that annihilates the holonomy (\xref{p09:sec:failures}). It is \emph{not} to fill the void with a shared image, for the idol forecloses the very openness in which real value is generated (\xref{p09:sec:non-possess}). It is \emph{not} to settle, by sign, a value not yet walked (\xref{p09:sec:failures}); \emph{not} to extract, hasten, or force-close (\xref{p09:sec:praxis}); \emph{not} to shape the manifold by the strong alone (\xref{p09:sec:self-legislation}). Cultivation is the positive name for the sum of these refusals---which is to say, it is the positive name for non-possession (\xref{p09:sec:non-possess}).

\begin{claim}[Cultivation is the tending of curvature, not the drawing of paths]
To cultivate a field is to tend its curvature so that the paths the participants walk of their own dynamics accrue a positive phase---and then to refrain from drawing those paths. It is the shaping of a condition, not the prescription of a content: the gardener does not pull the shoot upward but tends the soil, the light, the water, and lets the shoot grow of itself. Cultivation is thus a labour, and a demanding one, but a labour of the second kind---a tending of conditions, a subtraction of obstructions (\xref{p09:sec:repair})---and never a labour of the first kind, the production of the outcome directly. One cannot cultivate by forcing; the forcing is precisely what uncultivates.
\end{claim}

What, then, does the tending of curvature consist in, concretely? Three things, each a positive correlate of a refusal. \emc{It consists in keeping the void open}---declining to fill the other's irreducibility with an image, so that the encircling in which real value is generated remains possible; this is the positive practice of perception (\xref{p09:sec:perception}) sustained into a disposition. \emc{It consists in the patience that lets the path be walked}---declining the shortcut, the immediacy, the settling sign, so that the area a path encloses can actually be enclosed; this is the practical form of ``slowness is a formal necessary condition'' (\xref{p09:sec:failures}). \emc{And it consists in the holding-open of branching}---declining to prescribe where the other must arrive, so that their derivation can be genuinely their own; this is the interpersonal form of ``the grammar can be freely betrayed'' (\xref{p09:sec:non-possess}). To keep the void open, to let the path be walked, to hold branching open: this is the whole of cultivation, and it is non-possession made into a daily practice.

\subsection{The co-creation of value: generation along divergent paths}
\label{p09:sec:cocreation-proper}

If cultivation tends the field, co-creation is what happens within it---and the account of value (\xref{p09:sec:value}) tells us precisely what co-creation can and cannot be.

It \emph{cannot} be the joint production of a value transmitted, like two hands raising a single weight, for real value is not a conserved quantity transmitted but a phase generated along a path (\xref{p09:sec:value}). Two people do not co-create real value by pooling it; there is no pool. They co-create it by each walking their own path of encircling witness, within a shared field, such that the paths are not the same path. This is the decisive and easily missed point:

\begin{claim}[Co-creation is generation along divergent paths, not the merging of one]
The co-creation of real value is not two participants walking one path together but each walking their own, particular path---of their own irreducible dynamics (\xref{p09:sec:eudaimonia})---within a shared field whose curvature lets each path accrue a positive phase, and whose topology lets the phases reflux to both. Co-creation requires divergence, not merger: were the two paths one, there would be no encircling of an irreducible other, only the mutual confirmation of a shared image (the imaginary). The real value is generated precisely in the gap between the divergent paths---in the witness each gives the other's particularity, which is real exactly because it is not one's own. Co-creation is therefore the opposite of fusion: it is the generative holding-apart of two who witness, across the gap, what neither could generate alone.
\end{claim}

This is the constructive content of co-authorship (\xref{p09:sec:coupled}), and it shows why co-authorship is the right figure and merger the wrong one. Two authors of a poem do not write one line in unison; they write divergently, and the poem is the irreducible whole their divergence generates---a whole neither would have written, and neither owns. So with the co-creation of relational value: it is generated in the divergence, witnessed across the gap, owned by no one, and refluxing to each. The merged dyad of the imaginary co-creates nothing, for it has collapsed the gap in which alone real value is made; it can only circulate, ever more thinly, the single image it began with---which is why the merger that feels like the height of love is, in the value-theoretic terms of this paper, sterile, and why it exhausts (the rose again) rather than sustains.

\subsection{The field is jointly cultivated, or it is not a field}
\label{p09:sec:joint-cultivation}

One final turn closes the section and binds it back to justice. The field cannot be cultivated by one participant for the others, for a field cultivated unilaterally is not a field but a presented structure wearing a field's name---however gentle, it is still one party shaping the condition under which the other's dynamics must unfold, which is the shaping of the manifold by the strong alone (\xref{p09:sec:self-legislation}).

\begin{claim}[Joint cultivation, or counterfeit field]
A field is genuinely a field only if it is jointly cultivated---if each participant takes part in tending the curvature under which all unfold, and none is merely the recipient of a condition another has shaped. The benevolent unilateral cultivation---one party lovingly making space for the other---is a counterfeit field: it may feel like generosity, but it places one as the cultivator and the other as the cultivated, and the cultivated, however well-tended, is not a co-author but a tended thing. The asymmetry is the more insidious for being kindly; the gardener who most lovingly tends the shoot still stands outside it as its tender. A field, unlike a garden, has no gardener outside it: all who flourish in it are also, in flourishing, tending it for the rest.
\end{claim}

So the cultivation of the field, pursued to its ground, returns to the demand that opened the practical sequence and will govern its close: that no one be merely acted upon---neither reduced to fuel in the flow (\xref{p09:sec:coupled}), nor excluded from shaping the manifold (\xref{p09:sec:self-legislation}), nor, now, positioned as the cultivated rather than a co-cultivator of the field. Each of these is the same exclusion seen at a different scale, and the practice that answers all three is the same: the relentless conversion of the acted-upon into a co-author---of value, of the manifold, and of the field itself. \emb{To cultivate the field and to co-create value is, in the end, one act practised at every scale: to make of each participant a co-author, and of the relation a poem that no one tends from outside and no one owns.} What remains---the securing of this against real asymmetries of power, and the legislative practice by which it is held---is the work of the section to come.
